\section{Related Work}
\label{sec:related}

\subsection{Machine-Learning Diagnosis of Analog Circuits}
Data-driven fault diagnosis of analog circuits is an active field. Typical studies simulate hard
faults (open and short circuits) or soft faults (a component shifted by a given percentage) under
component tolerances, extract features from node voltages or from the frequency or time response,
and train a classifier. Dieste-Velasco has addressed hard faults in transistor amplifier stages
with a pattern-recognition neural network that uses a few accessible measurements, after grouping
the faults that those measurements cannot tell apart~\cite{dieste2021}; hard faults under
tolerances with a fuzzy inference system tuned by particle swarm optimization~\cite{dieste2024};
and soft faults in a Sallen-Key band-pass filter with seven classical classifiers on six voltages,
measured at two points and three frequencies, where the classes are the nominal circuit and a low
and a high deviation of each of seven components, reaching \SI{97.9}{\percent}
accuracy~\cite{dieste2025}. Deep models that learn features from raw responses are also
common~\cite{yang2020}, and explainable classifiers on frequency-response features have been
applied to parametric faults~\cite{rajpal2026}.

Three aspects of this literature are relevant here. First, the circuits: the Sallen-Key band-pass
filter and the four-op-amp biquad high-pass filter recur as
benchmarks~\cite{dieste2025,chen2025,zhang2026,srimani2022}, and Chen \emph{et
al.}~\cite{chen2025} show that their symmetries produce overlapping fault responses, which they
quantify with the distance between class centroids and the variance within classes over 100
training runs. Second, generalization across fault severities is recognized as an open problem:
Zhang \emph{et al.}~\cite{zhang2026} combine diffusion-based augmentation with domain adaptation
to transfer from a \SI{50}{\percent} to a \SI{25}{\percent} parametric deviation on those two
circuits and reach \SI{71}{\percent} accuracy on thirteen classes, a figure they present as an
upper bound. Third, the limits of what can be diagnosed are known from testability theory, which
defines ambiguity groups as sets of components whose faults cannot be distinguished from the
available measurements~\cite{starzyk2000}. In all of this work a fault is a deviation of a
component, usually by a fixed percentage of 25 or \SI{50}{\percent} against tolerances of 5 to
\SI{10}{\percent}~\cite{chen2025,zhang2026}; whether the circuit still meets its requirements is
not considered.

\subsection{Specification-Based and Alternate Test}
In production test, analog and mixed-signal circuits are checked against their specifications
because widely applicable fault models are lacking, and the cost of doing so has driven the
search for cheaper methods~\cite{cheng2010}. Alternate test predicts the specifications, or
directly the pass/fail outcome, from signatures that are cheap to measure. It rests on the
correlation between process parameters, signatures and specifications; the acceptance region in
signature space is obtained by Monte Carlo simulation and statistical classification, or a
regression from signatures to specifications is calibrated on a sample of measured
devices~\cite{cheng2010,voorakaranam2007}. Variyam and Chatterjee formulated test generation as
the reduction of the probability of misclassifying a bad circuit as good and vice
versa~\cite{variyam1998spec}, and addressed how the pass/fail threshold of alternate tests is
set~\cite{variyam1998synth}. Redundant specification tests can be removed by statistical
learning~\cite{biswas2005}. Machine-learning diagnosis is also presented at test venues, though on
the same benchmark filters and with a fault defined as a component outside its
tolerance~\cite{srimani2022}.

Alternate test has also been used for diagnosis. Cheng and Chang~\cite{cheng2010} review methods
in which the signatures, captured by built-in sensors, are related by statistical learning to the
source of a failure in integrated RF transceivers, and they call for test resources to be shared
with in-field testing. Our work follows the same principle and differs in its object and setting:
the components of a discrete network around an integrated amplifier, rather than the blocks of an
integrated circuit; a self-test run in service with the converter of the instrument, rather than
production or wafer-level test; the requirements of a clinical standard, rather than data-sheet
specifications; and a confounder, the electrode, that has no counterpart in integrated-circuit
test.

\subsection{Faults in ECG Acquisition and in Biomedical Equipment}
Work on faults in ECG acquisition operates on the recorded signal. FauDigPro~\cite{agarwal2022}
classifies ECG signals as normal or affected by a drift or a bias, both injected numerically into
recorded signals, and forecasts sensor failure; recent systems detect anomalies and sensor faults
in wearable and connected ECG devices~\cite{khezripour2026,adams2026}. These methods detect that
the data are bad. They do not say which part of the circuit has failed, nor whether the
instrument still meets its specifications.

At circuit level we have found two precedents in biomedical equipment. Gao~\cite{gao2025}
classifies faults of a simulated electroencephalograph circuit with wavelet-packet features and a
neural network and reports a missed-detection rate below \SI{9}{\percent}; the circuit, the
faults and their number are not described. Liu \emph{et al.}~\cite{liu2021} diagnose the control
board of a therapy device without schematics, from port signals and observed symptoms of seven
fault categories, with a long short-term memory network. Neither considers component tolerances,
compliance with specifications or the electrodes.

\subsection{The Skin-Electrode Interface}
The standard represents the skin-electrode impedance by \SI{51}{\kilo\ohm} in parallel with
\SI{47}{\nano\farad} in series with each lead~\cite{iec60601225}. Dry electrodes depart widely
from that figure. Joutsen \emph{et al.}~\cite{joutsen2024} fitted a single time-constant model to
impedance measured with dry electrodes of several materials and found parallel resistances from
about \SI{0.1}{\mega\ohm} for stainless steel to almost \SI{9}{\mega\ohm} for conductive fabric;
their open data show, for a given material, differences of one to three orders of magnitude
between subjects. An electrode impedance of megaohms forms a divider with the input network of a
conventional front-end and attenuates the signal, so a healthy instrument on a high-impedance
electrode can resemble a faulty one on a good electrode. Common-mode interference adds to this,
since any imbalance between electrodes converts part of it into a differential signal, which the
driven-right-leg circuit reduces but does not cancel~\cite{winter1983}. To our knowledge no
previous work separates electrode degradation from a fault of the circuit behind it.
